\cleardoublepage
\phantomsection
\section*{例题与作业题检索表}
\addcontentsline{toc}{section}{例题与作业题检索表}

本表按章节和题型汇总正文例题与作业题，页码由交叉引用自动生成，便于开卷考试时快速定位。

{\scriptsize
\renewcommand{\arraystretch}{1.18}

\noindent\textbf{第二章}\par
\begin{tabularx}{\textwidth}{@{}p{0.16\textwidth}p{0.13\textwidth}>{\raggedright\arraybackslash}X>{\raggedleft\arraybackslash}p{0.06\textwidth}@{}}
\toprule
\textbf{题型} & \textbf{题号} & \textbf{题目} & \textbf{页码} \\
\midrule
自信息与互信息 & 例~\ref{ex:ch2-ball-self-information} & 摸球问题：自信息、条件自信息与联合自信息 & \hyperref[ex:ch2-ball-self-information]{\pageref*{ex:ch2-ball-self-information}}  \\
自信息与互信息 & 例~\ref{ex:ch2-chessboard-information} & 棋盘问题：自信息与条件自信息 & \hyperref[ex:ch2-chessboard-information]{\pageref*{ex:ch2-chessboard-information}}  \\
熵与条件熵 & 例~\ref{ex:ch2-binary-source-entropy} & 二元信源熵及最大值 & \hyperref[ex:ch2-binary-source-entropy]{\pageref*{ex:ch2-binary-source-entropy}}  \\
熵与条件熵 & 例~\ref{ex:ch2-joint-entropy} & 联合熵、边缘熵与条件熵计算 & \hyperref[ex:ch2-joint-entropy]{\pageref*{ex:ch2-joint-entropy}}  \\
相对熵 & 例~\ref{ex:ch2-kl-divergence} & KL 散度计算 & \hyperref[ex:ch2-kl-divergence]{\pageref*{ex:ch2-kl-divergence}}  \\
条件熵 & 例~\ref{ex:ch2-three-boxes-conditional-entropy} & 三个盒子问题：条件熵比较 & \hyperref[ex:ch2-three-boxes-conditional-entropy]{\pageref*{ex:ch2-three-boxes-conditional-entropy}}  \\
熵的性质 & 例~\ref{ex:ch2-entropy-comparison} & 天气分布的熵比较 & \hyperref[ex:ch2-entropy-comparison]{\pageref*{ex:ch2-entropy-comparison}}  \\
平均互信息 & 例~\ref{ex:ch2-bsc-mutual-information} & 二元对称信道互信息 & \hyperref[ex:ch2-bsc-mutual-information]{\pageref*{ex:ch2-bsc-mutual-information}}  \\
平均互信息 & 例~\ref{ex:ch2-erasure-channel-mutual-information} & 二元删除信道互信息 & \hyperref[ex:ch2-erasure-channel-mutual-information]{\pageref*{ex:ch2-erasure-channel-mutual-information}}  \\
作业：信息量 & 作业~\ref{hw:ch2-bayes-information} & 贝叶斯公式与信息量 & \hyperref[hw:ch2-bayes-information]{\pageref*{hw:ch2-bayes-information}}  \\
作业：熵与互信息 & 作业~\ref{hw:ch2-entropy-conditional-mutual-information} & 熵、条件熵与互信息 & \hyperref[hw:ch2-entropy-conditional-mutual-information]{\pageref*{hw:ch2-entropy-conditional-mutual-information}}  \\
作业：平均互信息 & 作业~\ref{hw:ch2-weather-mutual-information} & 天气预报准确率与平均互信息 & \hyperref[hw:ch2-weather-mutual-information]{\pageref*{hw:ch2-weather-mutual-information}}  \\
\bottomrule
\end{tabularx}

\noindent\textbf{第三章}\par
\begin{tabularx}{\textwidth}{@{}p{0.16\textwidth}p{0.13\textwidth}>{\raggedright\arraybackslash}X>{\raggedleft\arraybackslash}p{0.06\textwidth}@{}}
\toprule
\textbf{题型} & \textbf{题号} & \textbf{题目} & \textbf{页码} \\
\midrule
无记忆信源 & 例~\ref{ex:binary-memoryless} & 二元无记忆信源熵与扩展熵 & \hyperref[ex:binary-memoryless]{\pageref*{ex:binary-memoryless}}  \\
无记忆信源 & 例~\ref{ex:ch3-second-extension-model} & 二次扩展模型 & \hyperref[ex:ch3-second-extension-model]{\pageref*{ex:ch3-second-extension-model}}  \\
无记忆信源 & 例~\ref{ex:ch3-binary-source-entropy} & 二元信源熵 & \hyperref[ex:ch3-binary-source-entropy]{\pageref*{ex:ch3-binary-source-entropy}}  \\
无记忆信源 & 例~\ref{ex:extension-entropy} & 二次扩展源熵 & \hyperref[ex:extension-entropy]{\pageref*{ex:extension-entropy}}  \\
平稳信源 & 例~\ref{ex:ch3-stationary-marginal} & 平稳信源的边缘分布 & \hyperref[ex:ch3-stationary-marginal]{\pageref*{ex:ch3-stationary-marginal}}  \\
马尔可夫链 & 例~\ref{ex:ch3-transition-matrix-check} & 判断转移矩阵 & \hyperref[ex:ch3-transition-matrix-check]{\pageref*{ex:ch3-transition-matrix-check}}  \\
马尔可夫链 & 例~\ref{ex:ch3-one-two-step-transition} & 一步、两步转移 & \hyperref[ex:ch3-one-two-step-transition]{\pageref*{ex:ch3-one-two-step-transition}}  \\
马尔可夫信源 & 例~\ref{ex:second-order} & 二阶马氏源建模 & \hyperref[ex:second-order]{\pageref*{ex:second-order}}  \\
马尔可夫信源 & 例~\ref{ex:markov-entropy} & 二元马氏链熵率与扩展熵 & \hyperref[ex:markov-entropy]{\pageref*{ex:markov-entropy}}  \\
作业：马氏链 & 作业~\ref{hw:ch3-homogeneous-markov-entropy} & 齐次马氏链熵计算 & \hyperref[hw:ch3-homogeneous-markov-entropy]{\pageref*{hw:ch3-homogeneous-markov-entropy}}  \\
作业：马氏源 & 作业~\ref{hw:ch3-symmetric-first-order-markov-source} & 对称一阶马氏源 & \hyperref[hw:ch3-symmetric-first-order-markov-source]{\pageref*{hw:ch3-symmetric-first-order-markov-source}}  \\
\bottomrule
\end{tabularx}

\noindent\textbf{第四章}\par
\begin{tabularx}{\textwidth}{@{}p{0.16\textwidth}p{0.13\textwidth}>{\raggedright\arraybackslash}X>{\raggedleft\arraybackslash}p{0.06\textwidth}@{}}
\toprule
\textbf{题型} & \textbf{题号} & \textbf{题目} & \textbf{页码} \\
\midrule
高斯差熵 & 例~\ref{ex:gauss-linear-transform} & 线性变换下的高斯熵 & \hyperref[ex:gauss-linear-transform]{\pageref*{ex:gauss-linear-transform}}  \\
高斯差熵 & 例~\ref{ex:gauss-orthogonal} & 正交变换下的高斯熵 & \hyperref[ex:gauss-orthogonal]{\pageref*{ex:gauss-orthogonal}}  \\
连续互信息 & 例~\ref{ex:gauss-mi} & 二维高斯互信息 & \hyperref[ex:gauss-mi]{\pageref*{ex:gauss-mi}}  \\
连续互信息 & 例~\ref{ex:gauss-signal-noise} & 高斯信号加噪声互信息 & \hyperref[ex:gauss-signal-noise]{\pageref*{ex:gauss-signal-noise}}  \\
作业：高斯变换 & 作业~\ref{hw:ch4-independent-gaussian-linear-combination} & 独立高斯变量的线性组合 & \hyperref[hw:ch4-independent-gaussian-linear-combination]{\pageref*{hw:ch4-independent-gaussian-linear-combination}}  \\
作业：高斯变换 & 作业~\ref{hw:ch4-correlated-gaussian-linear-transform} & 相关高斯变量与二维线性变换 & \hyperref[hw:ch4-correlated-gaussian-linear-transform]{\pageref*{hw:ch4-correlated-gaussian-linear-transform}}  \\
\bottomrule
\end{tabularx}

\noindent\textbf{第五章}\par
\begin{tabularx}{\textwidth}{@{}p{0.16\textwidth}p{0.13\textwidth}>{\raggedright\arraybackslash}X>{\raggedleft\arraybackslash}p{0.06\textwidth}@{}}
\toprule
\textbf{题型} & \textbf{题号} & \textbf{题目} & \textbf{页码} \\
\midrule
码的可译性 & 例~\ref{ex:ch5-unique-decodability} & 唯一可译性判断 & \hyperref[ex:ch5-unique-decodability]{\pageref*{ex:ch5-unique-decodability}}  \\
Kraft 不等式 & 例~\ref{ex:ch5-kraft-length-feasibility} & 码长可行性判断 & \hyperref[ex:ch5-kraft-length-feasibility]{\pageref*{ex:ch5-kraft-length-feasibility}}  \\
编码效率 & 例~\ref{ex:ch5-binary-extension-code-efficiency} & 二元扩展码的平均码长与效率 & \hyperref[ex:ch5-binary-extension-code-efficiency]{\pageref*{ex:ch5-binary-extension-code-efficiency}}  \\
Huffman 编码 & 例~\ref{ex:ch5-huffman-coding} & 二元 Huffman 编码 & \hyperref[ex:ch5-huffman-coding]{\pageref*{ex:ch5-huffman-coding}}  \\
作业：定长编码 & 作业~\ref{hw:ch5-binary-source-fixed-length-coding} & 无记忆二元信源等长编码 & \hyperref[hw:ch5-binary-source-fixed-length-coding]{\pageref*{hw:ch5-binary-source-fixed-length-coding}}  \\
作业：Huffman 编码 & 作业~\ref{hw:ch5-huffman-coding} & Huffman 编码 & \hyperref[hw:ch5-huffman-coding]{\pageref*{hw:ch5-huffman-coding}}  \\
\bottomrule
\end{tabularx}

\noindent\textbf{第六章}\par
\begin{tabularx}{\textwidth}{@{}p{0.16\textwidth}p{0.13\textwidth}>{\raggedright\arraybackslash}X>{\raggedleft\arraybackslash}p{0.06\textwidth}@{}}
\toprule
\textbf{题型} & \textbf{题号} & \textbf{题目} & \textbf{页码} \\
\midrule
信道模型 & 例~\ref{ex:ch6-bsc-model} & 二元对称信道（BSC） & \hyperref[ex:ch6-bsc-model]{\pageref*{ex:ch6-bsc-model}}  \\
信道模型 & 例~\ref{ex:ch6-erasure-channel-model} & 二元删除信道 & \hyperref[ex:ch6-erasure-channel-model]{\pageref*{ex:ch6-erasure-channel-model}}  \\
对称信道容量 & 例~\ref{ex:ch6-strongly-symmetric-check} & 判断强对称信道 & \hyperref[ex:ch6-strongly-symmetric-check]{\pageref*{ex:ch6-strongly-symmetric-check}}  \\
对称信道容量 & 例~\ref{ex:ch6-strongly-symmetric-capacity} & 强对称信道容量 & \hyperref[ex:ch6-strongly-symmetric-capacity]{\pageref*{ex:ch6-strongly-symmetric-capacity}}  \\
级联信道 & 例~\ref{ex:ch6-cascaded-bsc} & 两级 BSC 级联 & \hyperref[ex:ch6-cascaded-bsc]{\pageref*{ex:ch6-cascaded-bsc}}  \\
级联信道 & 例~\ref{ex:ch6-cascaded-bsc-capacity} & 两级 BSC 容量 & \hyperref[ex:ch6-cascaded-bsc-capacity]{\pageref*{ex:ch6-cascaded-bsc-capacity}}  \\
作业：信道容量 & 6.3 & 2 输入 3 输出信道容量 & \hyperref[hw:ch6-3-channel-capacity]{\pageref*{hw:ch6-3-channel-capacity}}  \\
作业：级联信道 & 6.14 & BSC 与删除型信道级联容量 & \hyperref[hw:ch6-14-cascaded-bsc-erasure-capacity]{\pageref*{hw:ch6-14-cascaded-bsc-erasure-capacity}}  \\
\bottomrule
\end{tabularx}

\noindent\textbf{第七章}\par
\begin{tabularx}{\textwidth}{@{}p{0.16\textwidth}p{0.13\textwidth}>{\raggedright\arraybackslash}X>{\raggedleft\arraybackslash}p{0.06\textwidth}@{}}
\toprule
\textbf{题型} & \textbf{题号} & \textbf{题目} & \textbf{页码} \\
\midrule
判决译码 & 例~\ref{ex:ch7-decision-rule-error-rate} & 判决规则与平均错误率 & \hyperref[ex:ch7-decision-rule-error-rate]{\pageref*{ex:ch7-decision-rule-error-rate}}  \\
MAP/ML 判决 & 例~\ref{ex:ch7-map-ml-decision} & MAP 与 ML 判决 & \hyperref[ex:ch7-map-ml-decision]{\pageref*{ex:ch7-map-ml-decision}}  \\
MAP/ML 判决 & 例~\ref{ex:ch7-nonuniform-map-ml} & 不等概输入下 MAP 与 ML 的差异 & \hyperref[ex:ch7-nonuniform-map-ml]{\pageref*{ex:ch7-nonuniform-map-ml}}  \\
纠错编码 & 例~\ref{ex:ch7-minimum-distance-correction} & 最小距离与纠错能力 & \hyperref[ex:ch7-minimum-distance-correction]{\pageref*{ex:ch7-minimum-distance-correction}}  \\
序列译码 & 例~\ref{ex:ch7-bsc-sequence-ml} & BSC 上的序列 ML 译码 & \hyperref[ex:ch7-bsc-sequence-ml]{\pageref*{ex:ch7-bsc-sequence-ml}}  \\
信道编码定理 & 例~\ref{ex:ch7-bsc-capacity-time} & BSC 容量与传输时间 & \hyperref[ex:ch7-bsc-capacity-time]{\pageref*{ex:ch7-bsc-capacity-time}}  \\
作业：传输速率 & 7.2 & 固定码长下的信息平均传输速率 & \hyperref[hw:ch7-2-fixed-code-length-transmission-rate]{\pageref*{hw:ch7-2-fixed-code-length-transmission-rate}}  \\
作业：MAP/ML 判决 & 7.4 & BSC 中 MAP 与 ML 判决 & \hyperref[hw:ch7-4-bsc-map-ml-decision]{\pageref*{hw:ch7-4-bsc-map-ml-decision}}  \\
作业：MAP 译码 & 7.14 & 非等概码字的 MAP 译码与可靠性比较 & \hyperref[hw:ch7-14-nonuniform-map-reliability]{\pageref*{hw:ch7-14-nonuniform-map-reliability}}  \\
\bottomrule
\end{tabularx}

\noindent\textbf{第八章}\par
\begin{tabularx}{\textwidth}{@{}p{0.16\textwidth}p{0.13\textwidth}>{\raggedright\arraybackslash}X>{\raggedleft\arraybackslash}p{0.06\textwidth}@{}}
\toprule
\textbf{题型} & \textbf{题号} & \textbf{题目} & \textbf{页码} \\
\midrule
加性噪声信道 & 例~\ref{ex:ch8-amplitude-limited-uniform-noise} & 幅度受限的均匀噪声信道 & \hyperref[ex:ch8-amplitude-limited-uniform-noise]{\pageref*{ex:ch8-amplitude-limited-uniform-noise}}  \\
注水分配 & 例~\ref{ex:ch8-two-parallel-gaussian-e6} & 两路并联高斯信道，$E=6$ & \hyperref[ex:ch8-two-parallel-gaussian-e6]{\pageref*{ex:ch8-two-parallel-gaussian-e6}}  \\
注水分配 & 例~\ref{ex:ch8-two-parallel-gaussian-e15} & 两路并联高斯信道，$E=15$ & \hyperref[ex:ch8-two-parallel-gaussian-e15]{\pageref*{ex:ch8-two-parallel-gaussian-e15}}  \\
AWGN 容量 & 例~\ref{ex:ch8-awgn-capacity} & AWGN 信道容量计算 & \hyperref[ex:ch8-awgn-capacity]{\pageref*{ex:ch8-awgn-capacity}}  \\
带宽与可靠传输 & 例~\ref{ex:ch8-reliable-transmission-bandwidth} & 可靠传输与带宽判断 & \hyperref[ex:ch8-reliable-transmission-bandwidth]{\pageref*{ex:ch8-reliable-transmission-bandwidth}}  \\
作业：加性噪声容量 & 8.3 & 限幅加性均匀噪声信道容量 & \hyperref[hw:ch8-3-amplitude-limited-uniform-noise-capacity]{\pageref*{hw:ch8-3-amplitude-limited-uniform-noise-capacity}}  \\
作业：注水分配 & 8.10 & 四路并联高斯信道注水分配 & \hyperref[hw:ch8-10-parallel-gaussian-water-filling]{\pageref*{hw:ch8-10-parallel-gaussian-water-filling}}  \\
作业：AWGN 容量 & 8.16 & AWGN 波形信道容量计算 & \hyperref[hw:ch8-16-awgn-waveform-capacity]{\pageref*{hw:ch8-16-awgn-waveform-capacity}}  \\
\bottomrule
\end{tabularx}

\noindent\textbf{第九章}\par
\begin{tabularx}{\textwidth}{@{}p{0.16\textwidth}p{0.13\textwidth}>{\raggedright\arraybackslash}X>{\raggedleft\arraybackslash}p{0.06\textwidth}@{}}
\toprule
\textbf{题型} & \textbf{题号} & \textbf{题目} & \textbf{页码} \\
\midrule
率失真端点 & 例~\ref{ex:ch9-dmin-dmax} & $D_{\min}$ 与 $D_{\max}$ 计算 & \hyperref[ex:ch9-dmin-dmax]{\pageref*{ex:ch9-dmin-dmax}}  \\
限失真编码 & 例~\ref{ex:ch9-single-symbol-compression} & 单符号压缩算法 & \hyperref[ex:ch9-single-symbol-compression]{\pageref*{ex:ch9-single-symbol-compression}}  \\
限失真传输 & 习题9.25 & 二元信源的限失真传输 & \hyperref[ex:ch9-binary-source-distortion-transmission]{\pageref*{ex:ch9-binary-source-distortion-transmission}}  \\
\bottomrule
\end{tabularx}
}
